\clearpage
\section*{Appendix X. The axisymmetric pressure-mode filter at a shrinking H8 interface}
\addcontentsline{toc}{section}{Appendix X. Axisymmetric pressure-mode filter}
\noindent\textbf{Scope.} We prove an exact cancellation for the leading \emph{carrier self-pressure} in the conditional transverse axisymmetric limit. This is not an authentic H8 calculation: the original source-dependent nonaxisymmetric shapers, the full periodic pressure of the nonlinear correction, and the source-matched positive stress are not yet fixed.

\paragraph{The reference geometry and leading solenoidality.}
Fix $0<h<1/6$, and let $D=1/2-h$, $A=1/2+h$, $p=1/2-3h$, and $\tau=T-t$. As in Appendices U and W, put $S_\tau=\operatorname{diag}(\tau^{1/2},\tau^{1/2},\tau^D)$ and let $\Omega(t)=x_0+S_\tau\Omega_0$ in a torus chart. Assume $\Omega_0$ is a smooth region of revolution whose horizontal sections are disks $\{R<a(\eta)\}$, with a regular smooth tip at either endpoint. Assume $U$ vanishes on an open collar of $\partial\Omega(t)$, and that the solenoidal exterior correction $w$ admits differentiated reference-coordinate profiles
\[
 \tau^{1/2}w_\perp(t,x_0+S_\tau y)\longrightarrow V_\perp(y),\qquad
 \tau^{1/2}w_z(t,x_0+S_\tau y)\longrightarrow V_z(y)
\]
in $C^1$ on each closed reference cross-section, uniformly away from the tips (with integrable uniform control at the tips). The physical incompressibility identity is
\[
 0=\operatorname{div}w
 =\tau^{-1}\operatorname{div}_{y_\perp}(\tau^{1/2}w_\perp)
 +\tau^{-1+h}\partial_\eta(\tau^{1/2}w_z).
\]
Thus $\operatorname{div}_{y_\perp}V_\perp=0$ on every reference disk and, by the planar divergence theorem,
\begin{equation}\label{eq:X-meanflux}
 \boxed{\int_0^{2\pi}q(\theta,\eta)\,d\theta=0,\qquad
 q(\theta,\eta)=V_\perp(a(\eta),\theta,\eta)\cdot e_R(\theta).}
\end{equation}

\paragraph{Lemma X.1 (Pressure generated by compact axisymmetric swirl in the transverse limit).}
Let $v_\theta(R,\eta)$ be an axisymmetric regular azimuthal velocity with $v_\theta=0$ for $R\ge R_*(\eta)<a(\eta)$. Assume the leading carrier self-pressure $\Pi_U$ is the regular radial solution, up to an additive function of $\eta$, of the transverse Poisson equation generated by $v_\theta e_\theta$:
\begin{equation}\label{eq:X-2d-poisson}
 -\Delta_{y_\perp}\Pi_U
 =\partial_{y_i}\partial_{y_j}
 [(v_\theta e_\theta)_i(v_\theta e_\theta)_j]
 =-\frac1R\partial_R(v_\theta^2).
\end{equation}
Then $\partial_R\Pi_U=v_\theta^2/R$ and $\Pi_U=C(\eta)$ in the entire exterior radial collar $R\ge R_*(\eta)$. In particular its trace on the reference boundary has no azimuthal Fourier modes except order zero.
\begin{proof}
The convective acceleration of axisymmetric pure swirl is $-(v_\theta^2/R)e_R$. Taking its transverse divergence proves the second equality in \eqref{eq:X-2d-poisson}. Radial integration yields $\partial_R\Pi_U=v_\theta^2/R+c(\eta)/R$; regularity at $R=0$ forces $c(\eta)=0$. Since the swirl vanishes on the collar, the radial pressure derivative vanishes there.
\end{proof}

\paragraph{Theorem X.2 (Vanishing carrier-only leading pressure work).}
Assume the hypotheses of Lemma X.1 and the leading-solenoidality conditions above. With the normal cofactor of Appendix W, the leading carrier self-pressure coefficient is
\begin{equation}\label{eq:X-carrier-zero}
 \boxed{C_{\pi,U}=\int_{\partial\Omega_0}\Pi_U(y)V_\perp(y)\cdot n_{0,\perp}(y)\,dS_0=0.}
\end{equation}
The claim is about the leading radial self-pressure $\Pi_U$, not the full pressure $\Pi$ of $U+w$.
\begin{proof}
Parametrize the lateral surface of revolution as $y=(a(\eta)\cos\theta,a(\eta)\sin\theta,\eta)$. Its oriented lateral surface element is $n_{0,\perp}dS_0=a(\eta)e_R(\theta)\,d\theta d\eta$. The pressure trace is $C(\eta)$ by Lemma X.1. Hence
\[
 C_{\pi,U}=\int a(\eta)C(\eta)\bigg[\int_0^{2\pi}q(\theta,\eta)\,d\theta\bigg]d\eta=0
\]
by \eqref{eq:X-meanflux}.
\end{proof}

\paragraph{Corollary X.3 (Exact angular-mode selection).}
Let $\Pi(\theta,\eta)$ be any sufficiently regular leading \emph{full} pressure on the boundary, and expand $\Pi(\theta,\eta)=\sum_{m\in\mathbb Z}\Pi_m(\eta)e^{im\theta}$ and $q(\theta,\eta)=\sum_{m\in\mathbb Z}q_m(\eta)e^{im\theta}$. Then
\begin{equation}\label{eq:X-angular-modes}
 \boxed{C_\pi=2\pi\int a(\eta)\sum_{m\ne0}\Pi_m(\eta)q_{-m}(\eta)\,d\eta.}
\end{equation}
The axisymmetric $m=0$ pressure channel cancels \emph{exactly}, because $q_0=0$ at each $\eta$. A successful pressure-only mechanism must supply a nonzero, signed overlap of matching $m\ne0$ pressure and normal-velocity modes. Such overlap may arise from the original nonaxisymmetric H8 shapers, the full $2\mathcal Q(U,w)$ interaction, or $\pi[w]$, but it has not been demonstrated for the authentic source. The other energy fluxes remain relevant and can also supply the needed inward transfer.

\paragraph{An exact prepared-field control.}
For $R<1$, take $v_\theta(R,\eta)=R(1-R^2)^3b(\eta)$ and extend it by zero for $R\ge1$; this is a finite-differentiability diagnostic, not an H8 source. A radial self-pressure is
\[
 \Pi_U(R,\eta)=-\tfrac1{14}(1-R^2)^7b(\eta)^2\quad(0\le R\le1),
 \qquad \Pi_U(R,\eta)=0\quad(R\ge1).
\]
Direct differentiation gives $\partial_R\Pi_U=v_\theta^2/R$. It is therefore identically constant on every exterior circular interface. An angular pressure mode can nevertheless have nonzero signed pairing: for $q=\cos\theta$ and $\Pi=\cos\theta$, one has $\int_0^{2\pi}\Pi q\,d\theta=\pi$, whereas $\Pi\equiv1$ gives zero. These are source-independent diagnostics.

\paragraph{Scope of the remaining physical theorem.}
The exact identity \eqref{eq:X-angular-modes} is a useful \emph{mode-selection gate}, not an existence proof. It presumes the stated differentiated anisotropic boundary limits; it does not derive the transverse pressure limit for the original nonaxisymmetric, rotating-frame H8 source. Before applying it, one must freeze the genuine Case~006 vector potential, common cutoff, all shapers and overlaps, rotated stress frame, density and stress centre, periodic Poisson pressure, and the exterior correction solving the full unforced nonlinear equation. The preceding nine algebraic checks validate only the model identities. The complete unforced Navier--Stokes theorem remains OPEN.
